\documentclass[12pt,a4paper]{article}
\usepackage[margin=1in]{geometry}
\usepackage[T1]{fontenc}
\usepackage{lmodern}
\usepackage{setspace}
\usepackage{enumitem}
\usepackage{hyperref}
\usepackage{titlesec}
\usepackage{parskip}

\setstretch{1.15}
\setlist[itemize]{leftmargin=1.5em}
\setlist[enumerate]{leftmargin=1.8em}

\begin{document}

\begin{center}
{\Large\bfseries DSCPP-III-2026-T3 PROJECT PROPOSAL\par}
\vspace{1em}
{\large\bfseries PROJECT AND TEAM INFORMATION\par}
\end{center}

\section*{Project Title}
\textit{Cipher Valley: A Gamified Environment for Learning Programming and Developing Computational Problem-Solving Skills}

\section*{Student / Team Information}

\textbf{Team Name:} Cipher Valley\\
\textbf{Team \#:} T334

\subsection*{Team Member 1 (Team Lead)}
\textbf{Name:} ABHISHAL SINGH\\
\textbf{Student ID:} 2510360145\\
\textbf{Email:} 2510360145@geu.ac.in

\subsection*{Team Member 2}
\textbf{Name:} NAVNEET RANA\\
\textbf{Student ID:} 2510012635\\
\textbf{Email:} 2510012635@geu.c.in

\subsection*{Team Member 3}
\textbf{Name:} BHAVSHAYA BHARDWAJ\\
\textbf{Student ID:} 2510012454\\
\textbf{Email:} 

\subsection*{Team Member 4}
\textbf{Name:} \\
\textbf{Student ID:} \\
\textbf{Email:} 

\newpage
\begin{center}
{\large\bfseries PROPOSAL DESCRIPTION (10 pts)\par}
\end{center}

\section*{Motivation (1 pt)}
In modern day programmers have potential and the analytical skills to learn but they get easily boared by the traditional way of textbook, smartboard and lecture learning and cant seem to keep up the motivation or the enthusiasm. they need something that keep them on their edge and give them a new way of learning.

This is where cipher valley comes in the picture it gives a new gamified approach to learning with quests, interactive story mode, and many more yet to be brainstromed.

The motivation for this project occurred when we observe that people mostly enjoy story driven games like GTA V, Minecraft and Stardew Valley. Nowadays and analyse it. Then it occurred to me ``what if we combine a story driven game with programming that's how cipher valley started.''

The project aims to try to explore whether combining programming practice with interactive progression, contextual hints, feedback, and increasing difficulty can provide a more engaging problem-solving experience for beginner programmers.

\section*{State of the Art / Current solution (1 pt)}
\begin{itemize}
\item Gamification is also being introduced in many other ways to overcome this obstacle between the programmer and the joy of coding. The games which utilizes this gamification approach are:
\item Minecraft Education Edition: it helps student with basic codes and other subjects.
\item Baba is You: a 10/10 rated game on Steam which forces the student to think out of box.
\item Prodigy Math: where millions of users complete maths designed quests.
\item The other ways than gamification that are use to solve this problem are listed below:
\end{itemize}

\begin{enumerate}
\item interactive lectures
\item practical based approach
\item Relating the topics with real world examples
\item Quiz
\item group discussion and many more methods like competition.
\end{enumerate}

Cipher Valley uses these approaches and methods and tries to be a project that could also produce similar result or outcome using the gamification approach and most of all make programming fun and interesting.

\section*{Project Goals and Milestones (2 pts)}
The primary goal of Cipher Valley is to design and develop a functional prototype of a gamified programming-learning environment and investigate its potential to improve learner engagement and problem-solving experience.

\begin{itemize}
\item \textbf{Milestone 1 --- Research and Requirements:} Review relevant literature, identify the problem, define research questions, and establish educational and technical requirements.
\item \textbf{Milestone 2 --- System and Game Design:} learning to combine game and programming education to get the best outcome.
\item \textbf{Milestone 3 --- designing the story base:} trying to make a story which is interactive and keep people engaged and not feel bored.
\item \textbf{Milestone 4 --- Learning System \& C++/DSA Design ---} Decide which programming concepts, challenges, C++ components, graphs, BFS/DFS/Dijkstra, etc. fit into the game and where they will be used.
\item \textbf{Milestone 5 --- prototype development:} Making a working prototype of the game with the help of all the data, research and information acquired.
\item \textbf{Milestone 6 --- Testing and Evaluation:} Test the prototype with users and collect measures such as challenges attempted and completed, attempts, hint usage, completion time, and structured feedback.
\item \textbf{Milestone 7 --- Analysis and Documentation:} Analyse results, document limitations and findings, and prepare the final report, presentation, demonstration, and GitHub documentation.
\end{itemize}

\newpage
\section*{Project Approach (3 pts)}
The project will follow an iterative research-and-development approach. The first stage will involve reviewing literature on gamification, game-based learning, and programming education to establish the research foundation. Requirements will then be defined for both educational and technical components.

The prototype will be designed as a browser-based 2D interactive environment. Phaser.js and JavaScript will be used for the game layer, with HTML/CSS supporting the interface. Initial learning content will focus on variables, conditional statements, loops, functions, and basic debugging.

Programming challenges will be integrated into quests and game progression. A challenge engine will present problems and process submissions, while a hint and feedback mechanism will provide guidance without immediately revealing solutions. A progression system will record completed challenges and unlock subsequent content.

C++ will be incorporated for algorithmic and Data Structures and Algorithms components. Graph-based representations can be used for world navigation, with BFS, DFS, and Dijkstra's algorithm applied where appropriate. JSON or SQLite may be used for structured game and progress data. Git and GitHub will be used for version control and documentation.

After implementation, user testing will collect quantitative interaction data and structured feedback. Results will be analysed to evaluate engagement and problem-solving experience and to identify potential benefits and limitations.

\newpage
\section*{System Architecture (High Level Diagram)(2 pts)}
% The source Word document contains the heading/instruction for the diagram
% but no architecture diagram content. Insert the diagram here when available.

\section*{Project Outcome / Deliverables (1 pts)}
\begin{itemize}
\item The primary outcome will be a functional prototype of Cipher Valley demonstrating the integration of programming education with a game-based environment. The prototype is expected to include an interactive 2D game environment, programming challenges, quests or progression, contextual hints, immediate feedback, and basic progress tracking.
\item Learning the implementation of the DSA, C++ with the help of this project and gaining a better understanding of the concepts and its implementation in real world
\item And also game implementation of D.S.A and C++.
\item Deliverables include a maintained git hub repository link, a well documented findings as well as a full understanding of implementation of game development experienced gain and knowledge learned in a file or document formant.
\item And most importantly what we learned through this PBL
\end{itemize}

\section*{Assumptions}
The project assumes that initial target users have basic computer literacy and are beginners or early-stage programming learners. Users are assumed to have access to a computer with a modern web browser. The target has a prior interaction with games or enjoy playing them. The results of the prototype will be considered as a bench mark rather than real world application.

\section*{References}
\begin{enumerate}
\item University Refrenced books for DSA integration and C++ programming, which is suggested by university for learning the basic and fundamentals as well as advanced topics of the subject.
\item Gamefied learning in higher education by Subash and Cundey (2018):\\
\url{https://www.sciencedirect.com/science/article/abs/pii/S0747563218302541}
\item The effectiveness of game introduction in programming by Guo and Lan:\\
\url{https://www.researchgate.net/publication/379980212_Investigating_the_influence_of_gamification_on_motivati\%0Aon_and_learning_outcomes_in_online_language_learning}
\end{enumerate}

\end{document}
